\chapter*{Jargon}
\addcontentsline{toc}{chapter}{jargon} 

\openepigraph{%
	There was in Italy a hidden demand for a boring government which would try to tell the truth in non-political jargon.
}{Mario Monti}

I hate jargon, even if I use it all the time. Jargon exists to 'simplify' complex exchanges amongst people who share an intimate understanding of some particular area of knowledge. Nuclear engineers have their jargon. Lobster fisherman have their jargon. Police have their jargon. Jargon may also vary geographically. Its use does the opposite of simplify when 'outsiders' attempt to participate in the conversation.

Systems thinkers are no exception to this rule. The many tribes of systems thinking each have contribute their own jargon, much of which is uninterpretable to the outsider. Jargon collapsed nuanced meaning into a word or phrase that is lost to those outside the tribe.

A road block on our systems thinking journey. A barrier to an entire paradigm of thinking that we need in order to increase the effectiveness of our public health practice. So let's begin a list of key jargon:

\newthought{complexity}. Just try and find agreement amongst systems thinkers around the meaning of this fundamental idea! I dare you. Snowden has a neat little diagram, Jackson references 10 common elements... What everyone is trying to say is that whatever complexity is, and it is many things, it means that wrapping our heads around the issues to which it is being applied is going to be non-trivial. We need to pause and think about this.

\newthought{system}. You know. Like 'health system', 'computer system', 'earth system', 'political system', but so is 'an organisation', 'a society', 'a hockey team', 'a car' or even 'a fuel pump' inside a car. While there might not be a lot of debate about what a system is, its still jargon to the outsider. A system is a \emph{perceived} entity made up of parts and having properties that none of its parts have on their own. Its 'a perceived entity', because unlike a car or computer, you can't touch or see many of the systems we are interested in. How you 'see' an organisation or a 'health system' is usually through someone's description (maybe your own) of the system.

\newthought{systems thinking}. A way of thinking about public health challenges. We conceive of these challenge as entities being a participant of something bigger that we call a 'system'. Our public health challenges are participants within a large system. Our challenge itself may be a system, and considered a sub-system of the larger system. Let's take obesity as an example. Obesity is a categorisation of human body weight that suggests an excess of body fat. In turn, obesity is a state of a system called the 'human body' or prosaicly 'a person', which most relevantly is composed of a systems that store fat, regulate hormones, put nutrients and calories that we eat to work, physical activity undertaken, etc. However, the systems perspective of obesity would be to look at the system of which obesity is just one part...

\newthought{thinking}.

\newthought{problems}. These are what you are given in school mathematics class. They probably also provide you with a definitive answer in the back of your textbook. If you like puzzles like a good crossword or something physical like a Rubic's cube, then you know about another type of problem. These too have definitive solutions. Problems have solutions. You apply the solution, once you know it, and your problem is solved. However, beyond puzzles and maths class, we don't deal with too many problems in public health, or more correctly, we have problems without definitive solutions. A different class of problem. One that might be called 'messy' or 'complex'. So to avoid confusion and differentiate between problems with known and definitive solutions and the real world we work in, let's agree to drop the term 'problem' and instead run with 'problematic situation'.

\newthought{problematic situation}

\newthought{whole system approach}. Seems quite popular in public health academic circles, but it doesn't feel right to me. I know what they mean, I think, in urging colleagues to look at everything connected to the problematic situation, not just the situation of interest itself. However, this language runs counter to the useful application of systems thinking, because (in the words of Donella Meadows) "there are no separate systems. The world is a continuum. Where to draw a boundary around a system depends on the purpose of the discussion." Let's avoid this term, it pushes a mindset that is at odds with the adoption of a cohesive systems practice. We never look at the whole system. Ever.

